\chapter{Prerequisiti}\stepcounter{capitoloeccolo}

\section{Prerequisiti di Logica}

\subsection{Logica del prim'ordine}

\definizione{
    Un \emph{linguaggio del prim'ordine} $ \mathcal{L} $ è un insieme costituito da \begin{itemize}
        \item un insieme numerabile di variabili $ \set{v_1,\dots,v_{n},\dots } $;
        \item un insieme di simboli logici, ovvero i connettivi e i quantificatori; 
        \item il simbolo di uguaglianza $ \eqcirc $;
        \item un insieme $ \set{R_{i}: i \in I } $ di simboli di relazione; 
        \item un insieme $ \set{F_{j}: j \in J } $ di simboli di funzione; 
        \item una funzione di arietà \begin{equation*}
            \operatorname{ar}: I \cup J \longrightarrow \N\setminus\set{0};
        \end{equation*}
        \item un insieme $ \set{c_{k}: k \in K } $ di simboli di costanti.
    \end{itemize}
}
\definizione{
    Da un linguaggio del prim'ordine $ \mathcal{L} $, una $ \mathcal{L} $-struttura $ \mathcal{A} $ è il dato di \begin{itemize}
        \item un insieme non vuoto $ A $, detto \emph{universo} o dominio (denotato con $ \norma{\mathcal{A}} $); 
        \item per ogni $ i \in I $, un insieme $ R_{i}^{\mathcal{A}} \subseteq A^{\operatorname{ar}(i)}  $;
        \item per ogni $ j \in J $, una funzione \begin{equation*}
            f^{\mathcal{A}}_j : A^{\operatorname{ar}(j)} \longrightarrow A;
        \end{equation*}
        \item per ogni $ k \in K $ un elemento $ c_{k}^{\mathcal{A}} \in A  $.
    \end{itemize}
}

\definizione{
    Date due $ \mathcal{L} $-strutture $ \mathcal{A}, \mathcal{B} $, un morfismo \begin{equation*}
        F: \mathcal{A} \longrightarrow \mathcal{B}
    \end{equation*}è una mappa $ \pi: \norma{\mathcal{A}} \longrightarrow \norma{ \mathcal{B}} $ che soddisfi \begin{itemize}
        \item[(A)] per ogni $ i \in I $: \begin{equation*}
            \left(a_1,\dots,a_{\operatorname{ar}(i)} \right) \in R_{i}^{\mathcal{A}} \,\implies\,\left(F\left(a_1\right),\dots,F\left(a_{\operatorname{ar}(i)} \right)\right) \in R_{i}^{ \mathcal{B}}; 
        \end{equation*}
        \item[(B)] per ogni $ j \in J $: \begin{equation*}
            F\left(f_{j}^{\mathcal{A}}\left(a_1,\dots,a_{\operatorname{ar}(j)} \right) \right) = f_{j}^{ \mathcal{B}}\left(F(a_1),\dots,F(a_{\operatorname{ar}(j)} )\right);
        \end{equation*}
        \item[(C)] per ogni $ k \in K $: $ F\left(c_{k}^{\mathcal{A}} \right) = c_{k}^{ \mathcal{B}}  $.
    \end{itemize}Inoltre $ F $ si dice \emph{morfismo pieno} se vale la condizione (A'): per ogni $ i \in I $\begin{equation*}
        \left(a_1,\dots,a_{\operatorname{ar}(i)} \right) \in R_{i}^{\mathcal{A}} \,\iff\,\left(F\left(a_1\right),\dots,F\left(a_{\operatorname{ar}(i)} \right)\right) \in R_{i}^{ \mathcal{B}}; 
    \end{equation*}Un morfismo pieno e iniettivo si dice \emph{immersione}.
}
\definizione{
    Siano $ \mathcal{A}, \mathcal{B} $ due $ \mathcal{L} $-strutture. $ \mathcal{A} $ è una sottostruttura di $ \mathcal{B} $ se \begin{itemize}
        \item $ \norma{\mathcal{A}} \subseteq \norma{ \mathcal{B}} $;
        \item $ \operatorname{id}: \mathcal{A} \longrightarrow \mathcal{B} $ è una immersione. 
    \end{itemize}Si scrive $ \mathcal{A} \subseteq \mathcal{B} $.
}

\definizione{
    Sono \emph{termini}: \begin{itemize}
        \item costanti e variabili;
        \item se $ f $ è un simbolo di funzione con $ \operatorname{ar} (f)=n $ e $ t_1,\dots,t_{n} $ sono termini, allora è termine:\[
            f(t_1,\dots,t_{n} )
        \]
    \end{itemize}
}
\osservazione{
    È naturale introdurre un concetto di \emph{altezza} di una termine: nel grafo costruito a partire termine $ t $, che ricostruisce le operazioni utilizzate per costruirla, l'altezza $ \altezza(t) $ è la lunghezza del ramo più lungo (partendo da zero).
}
\definizione{
Sia $ t(x_1,\dots,x_{n} ) $ un termine, dove $ x_1,\dots,x_{n}  $ sono le variabili che compaiono in $ t $, e dati $ s_1,\dots,s_{n}  $ termini, la scrittura: \[
    t\left[s_1/x_1,\,s_2/x_2,\dots,\,s_{n}/x_{n}  \right]
\]è il termine ottenuto da $ t $ sostituendo ogni $ x_{i}  $ con $ s_{i}  $.
}

\paragrafo{Lemma}{%
    $ t\left[s_1/x_1,\,s_2/x_2,\dots,\,s_{n}/x_{n}  \right] $ è un termine.
}{}{}
\definizione{
    Una \emph{formula atomica} è \begin{itemize}
        \item $ t_1 \eqcirc t_2 $ con $ t_1,t_2 $ termini; 
        \item se $ R_{i}  $ è una relazione di arietà $ n $, è formula atomica \begin{equation*}
            R_{i}(t_1,\dots,t_{n} ) 
        \end{equation*}per $ t_1,\dots,t_{n}  $ termini.
    \end{itemize}

    L'insieme delle \emph{formule} è definito induttivamente come la chiusura dell'insieme delle formule atomiche rispetto alle seguenti operazioni \begin{equation*}
        \varphi \mapsto \lnot\,\varphi,\quad \varphi,\psi \mapsto \varphi \,\lor\, \psi,\quad \varphi\mapsto \exists\, x \varphi.
    \end{equation*}

    Tutte le variabili che all'interno di una formula non sono sotto il raggio di un quantificatore si dicono \emph{libere}: indicando con $ \varphi(x_1,\dots,x_{n} ) $ si intende che le variabili libere di $ \varphi $ sono un sottoinsieme di $ \set{x_1,\dots,x_{n} } $.
}
\osservazione{
    È naturale introdurre un concetto di \emph{altezza} di una formula: nel grafo costruito a partire dalla formula $ \varphi $, che ricostruisce le operazioni utilizzate per costruirla, l'altezza $ \altezza(\varphi) $ è la lunghezza del ramo più lungo (partendo da zero).
}

\definizione{
    Sia $ t $ un $ \mathcal{L} $-termine e $ x_1,\dots,x_{n}  $ le variabili che occorrono in $ t $. Sia $ \mathcal{M} $ una $ \mathcal{L} $-struttura e $ a_1,\dots,a_{n} \in M  $. Allora \[
        t^{\mathcal{M}}[a_1,\dots,a_{n} ]
    \]è l'elemento di $ M $ definito per ricorsione da \begin{itemize}
        \item se $ t $ è nella forma $ x_{i} $, allora $ t^{\mathcal{M}}[a_1,\dots,a_{n} ] = a_{i} $;
        \item se $ t $ è nella forma $ c $ per $ c $ simbolo di costante, allora $ t^{\mathcal{M}} [a_1,\dots,a_{n} ] = c^{\mathcal{M}} $;
        \item se $ t $ è nella forma $ f(s_1,\dots,s_{k} ) $ per qualche $ f $ simbolo di funzione con arietà $ k $ e $ s_1,\dots,s_{k}  $ termini, allora \[
            t^{\mathcal{M}}[a_1,\dots,a_{n} ] = f^{\mathcal{M}}\left(s_1^{\mathcal{M}}[a_1,\dots,a_{n} ],\dots,s_k^{\mathcal{M}}[a_1,\dots,a_{n} ]\right)
        \]  
    \end{itemize}
}
\definizione{
    Data una $ \mathcal{L} $-struttura $ \mathcal{M} $ ed una formula $ \varphi(x_1,\dots,x_{n} ) $, si definisce la relazione di soddisfazione \[
        \mathcal{M}\vDash\varphi[a_1,\dots,a_{n} ]
    \]per ricorcorsione sulla complessità di $ \varphi(x_1,\dots,x_{n} ) $. \begin{itemize}
        \item Se $ \varphi(x_1,\dots,x_{n} ) $ è del tipo \[\varphi(x_1,\dots,x_{n} ):\hspace{4em} t\eqcirc s\](con $ t $ ed $ s $ termini), allora $ \mathcal{M}\vDash \varphi[a_1,\dots,a_{n} ] $ se e solo se \[
            t^{\mathcal{M}}[a_1,\dots,a_{n} ] = s^{\mathcal{M}}[a_1,\dots,a_{n} ].
        \]
        \item Se $ \varphi(x_1,\dots,x_{n} ) $ è del tipo \[\varphi(x_1,\dots,x_{n} ):\hspace{4em} R(t_1,\dots,t_{k} )\] per $ R $ relazione e $ t_{i}  $ termini, allora $ \mathcal{M}\vDash \varphi[a_1,\dots,a_{n} ] $ se e solo se \[
            \left(t_1^{\mathcal{M}}[a_1,\dots,a_{n} ],\dots,t_k^{\mathcal{M}}[a_1,\dots,a_{n} ]\right) \in R^{\mathcal{M}}.
        \]
        \item Se $ \varphi(x_1,\dots,x_{n} ) $ è del tipo \[\varphi(x_1,\dots,x_{n} ):\hspace{4em} \lnot\,\psi\]allora $ \mathcal{M}\vDash\varphi[a_1,\dots,a_{n} ] $ se e solo se \[
            \mathcal{M}\nvDash\psi[a_1,\dots,a_{n} ].
        \]
        % \item Se $ \varphi(x_1,\dots,x_{n} ) $ è del tipo \[%
        % \psi_1 \,\land\, \psi_2
        % \]allora $ \mathcal{M}\vDash\varphi[a_1,\dots,a_{n} ] $ se e solo se \[
        %     \mathcal{M}\vDash\psi_1[a_1,\dots,a_{n} ]\quad\text{e}\quad \mathcal{M}\vDash\psi_2[a_1,\dots,a_{n} ].
        % \]
        \item Se $ \varphi(x_1,\dots,x_{n} ) $ è del tipo \[%
        \psi_1 \,\lor\, \psi_2
        \]allora $ \mathcal{M}\vDash\varphi[a_1,\dots,a_{n} ] $ se e solo se \[
            \mathcal{M}\vDash\psi_1[a_1,\dots,a_{n} ]\quad\text{oppure (inclusivo)}\quad \mathcal{M}\vDash\psi_2[a_1,\dots,a_{n} ].
        \]
        % \item Se $ \varphi(x_1,\dots,x_{n} ) $ è del tipo \[%
        % \psi_1 \,\implies\, \psi_2
        % \]allora $ \mathcal{M}\vDash\varphi[a_1,\dots,a_{n} ] $ se e solo se \[
        %     \mathcal{M}\nvDash\psi_1[a_1,\dots,a_{n} ]\quad\text{oppure }\quad \mathcal{M}\vDash\psi_2[a_1,\dots,a_{n} ].
        % \]
        % \item Se $ \varphi(x_1,\dots,x_{n} ) $ è del tipo \[%
        % \psi_1 \,\iff\, \psi_2
        % \]allora $ \mathcal{M}\vDash\varphi[a_1,\dots,a_{n} ] $ se e solo se \[
        %     \mathcal{M}\vDash\psi_1[a_1,\dots,a_{n} ]\quad\text{se e solo se}\quad \mathcal{M}\vDash\psi_2[a_1,\dots,a_{n} ].
        % \]
        \item Se $ \varphi(x_1,\dots,x_{n} ) $ è del tipo $ \exists\, y\,\psi $ è necessario distinguere tre casi. \begin{itemize}
            \item Se $ y $ non occorre libera in $ \psi $, allora si può scrivere $ \psi $ come $ \psi(x_1,\dots,x_{n} ) $ e $ \mathcal{M}\vDash\psi[a_1,\dots,a_{n} ] $ se e solo se \[
                \mathcal{M}\vDash \psi[a_1,\dots,a_{n} ].
            \]
            \item Se $ y $ occorre libera in $ \psi $ ed è diversa da tutte le $ x_{i}  $, allora si può scrivere $ \psi $ come $ \psi(x_1,\dots,x_{n},y ) $, e $ \mathcal{M}\vDash\varphi[a_1,\dots,a_{n} ] $ se e solo se esiste qualche $ b \in M $ tale che \[
                \mathcal{M}\vDash \psi[a_1,\dots,a_{n},b ].
            \]
            \item Se $ y $ occorre libera in $ \psi $ è $ y=x_{i}  $ per qualche $ i $, allora si può scrivere $ \psi $ come $ \psi(x_1,\dots,x_{i-1},y,x_{i+1},\dots,x_{n}   ) $ e si avrà che $ \mathcal{M}\vDash\varphi[a_1,\dots,a_{m} ] $ se e solo se esiste qualche $  b \in M $ tale che \[
                \mathcal{M}\vDash\psi[a_1,\dots,a_{i-1},b,a_{i+1},\dots,a_{n}].
            \]
        \end{itemize}
        % \item Se $ \varphi(x_1,\dots,x_{n} ) $ è del tipo $ \forall\, y\,\psi $ si devono distinguere tre casi. \begin{itemize}
        %     \item Se $ y $ non occorre libera in $ \psi $, allora possiamo scrivere $ \psi $ come $ \psi(x_1,\dots,x_{n} ) $ e $ \mathcal{M}\vDash\psi[a_1,\dots,a_{n} ] $ se e solo se \[
        %         \mathcal{M}\vDash \psi[a_1,\dots,a_{n} ].
        %     \]
        %     \item Se $ y $ occorre libera in $ \psi $ ed è diversa da tutte le $ x_{i}  $, allora possiamo scrivere $ \psi $ come $ \psi(x_1,\dots,x_{n},y ) $, e $ \mathcal{M}\vDash\varphi[a_1,\dots,a_{n} ] $ se e solo se per ogni $ b \in M $ si ha che \[
        %         \mathcal{M}\vDash \psi[a_1,\dots,a_{n},b ].
        %     \]
        %     \item Se $ y $ occorre libera in $ \psi $ è $ y=x_{i}  $ per qualche $ i $, allora possiamo scrivere $ \psi $ come $ \psi(x_1,\dots,x_{i-1},y,x_{i+1},\dots,x_{n}   ) $ e si avrò che $ \mathcal{M}\vDash\varphi[a_1,\dots,a_{m} ] $ se e solo se per ogni $  b \in M $ si ha che \[
        %         \mathcal{M}\vDash\psi[a_1,\dots,a_{i-1},b,a_{i+1},\dots,a_{n}].
        %     \]
        % \end{itemize}
    \end{itemize}
}

\definizione{Siano $ \mathcal{N} \subseteq \mathcal{M} $ due $ \mathcal{L} $-strutture. $ \mathcal{N} $ è una \emph{sottostruttura elementare di $ \mathcal{M} $} se per ogni $ \mathcal{L} $-formula $ \varphi(x_1,\dots,x_{n} ) $ ed ogni $ a_1,\dots,a_{n} \in N$ \[
    \mathcal{M}\vDash\varphi[a_1,\dots,a_{n} ]\quad\text{se e solo se}\quad \mathcal{N}\vDash\varphi[a_1,\dots,a_{n} ],
\]}
\definizione{
    Se $ F:\mathcal{N}\longrightarrow \mathcal{M}$ è una immersione di $ \mathcal{N} $ in $ \mathcal{M} $ e $ F(\norma{\mathcal{N}}) $, visto come sottostruttura di $ \mathcal{M} $, è una sottostruttura elementare di $ \mathcal{M} $, si dice che $ F $ è una \emph{immersione elementare} e che $ \mathcal{N} $ si \emph{immerge elementarmente} in $ \mathcal{M} $, in simboli $ \mathcal{N}\preceq \mathcal{M} $.
}

Si riporta inoltre un lemma necessario per dimostrare in seguito un teorema.

\paragrafo{Lemma}{%
Sia $ F:\mathcal{N}\longrightarrow \mathcal{M} $ un'immersione. Le seguenti affermazioni sono equivalenti: \begin{enumerate}
    \item $ F $ è un'immersione elementare; 
    \item per ogni $ \mathcal{L} $-formula $ \varphi(y,x_1,\dots,x_{n} ) $ ed ogni $ a_1,\dots,a_{n} \in N  $ \begin{gather*}
        \mathcal{M}\vDash \exists y \, \varphi\left[y, F(a_1),\dots,F(a_{n}) \right]\\[1em] \text{se e solo se}\\[1em] \mathcal{M}\vDash\varphi\left[F(a_{0}), F(a_1),\dots,F(a_{n}) \right]\text{ per qualche }a_0 \in N.
    \end{gather*}
\end{enumerate}
}{odijcoijodijocijoijoijoi}{}

\subsection{Assioma della scelta}

Per la costruzione di $ \starR $ come ultraprodotto, sarà necessario utilizzare l'Assioma della Scelta. Lo si presenta anche nella sua forma equivalente: il Lemma di Zorn

\paragrafo{Assioma di scelta}{%
    Per ogni insieme $ X\neq \emptyset $ esiste \[f:\parti{X}\setminus\{\emptyset\}\longrightarrow X\]tale che $ f(Y) \in Y $ per ogni $ Y \subseteq X $ non vuoto.
}{}{}
\paragrafo{Lemma di Zorn}{%
    Se un insieme (parzialmente) ordinato è tale che ogni catena ha un maggiorante, allora ha un elemento massimale.
}{}{}


\section{Prerequisiti di Geometria Differenziale}\label{sec:geodiff}

\subsection{Varietà differenzabili}

\definizione{
    Una \emph{varietà topologica} $ n $-dimensionale è uno spazio topologico $ M $, di Hausdorf e a base numerabile, tale per cui, per ogni $ p \in M $ sono $ {U} $ intorno aperto di $ p $ in M, $ V \subseteq \R^{n} $ aperto e \begin{equation*}
        \phi : {U} \longrightarrow V
    \end{equation*}omeomorfismo. $ \phi $ si dice \emph{carta topologica}.
}
\definizione{
    Se $ M $ è una varietà topologica, un \emph{atlante topologico} di $ M $ è una famiglia di \emph{carte} \begin{equation*}
        \mathcal{A} = \set{(U_{\alpha}, \varphi_{\alpha}  )}_{\alpha \in I}
    \end{equation*}indicizzata su un insieme non vuoto, dove \begin{itemize}
        \item $ \set{U_{\alpha} }_{\alpha \in I} $ è un ricoprimento aperto di $ M $;
        \item $ \varphi_{\alpha}: U_{\alpha} \longrightarrow V_{\alpha}    $ è una carta topologica per certi $ V_{\alpha} \subseteq \R^{n}  $ aperti.
    \end{itemize}
}
\definizione{
    Sia $ M $ una varietà topologica e sia $ \mathcal{A} = \set{(U_{\alpha},\varphi_{\alpha}  )}_{\alpha \in I} $ un atlante topologico. $ \mathcal{A} $ è \emph{differenziabile} se tutti i cambi di carte sono dei diffeomorfismi tra aperti di $ \R^{n} $, ovvero se per ogni $ \alpha,\beta \in I $ tali che $ U_{\alpha}\cap U_{\beta}\neq \emptyset   $ \begin{equation*}
        \varphi_{\alpha}\circ \left(\varphi_{\beta}\right)^{-1}: \varphi_{\beta} (U_{\alpha} \cap U_{\beta}  )   \longrightarrow \varphi_{\alpha} (U_{\alpha} \cap U_{\beta}  )
    \end{equation*}è una biiezione continua e differenziabile (di classe $ C^{\infty} $), con inversa differenziabile (di classe $ C^{\infty} $).
}
\definizione{
    Due atlanti differenziabili si dicono \emph{compatibili} se la loro unione è ancora un atlante differenziabile.
}
\proposizione{kjndjckjncdjckjnsjnkjnkjn}{
    La compatibilità di atlanti è una relazione di equivalenza.
}
\definizione{
    Una \emph{struttura differenzabile} è un atlante differenziabile che sia massimale rispetto all'inclusione. 

    Una varietà topologica con una struttura differenziabile fissata si dice \emph{varietà differenziabile}.
}
\teorema{dkjnckskjncjskjndjckjndjdjjjj}{
    Ogni atlante differenziabile è contenuto in un'unica struttura differenziabile.
}

\subsection{Funzioni differenziabili tra varietà}

\definizione{
    Siano $ {M} $, $ {N} $ due varietà differenziabili, e siano $ \mathcal{A}=\set{(U_{\alpha}, \varphi_{\alpha}  )}_{\alpha \in I} $, $ \mathcal{B}=\set{(V_{\beta}, \psi_{\beta}  )}_{\beta \in J} $ le rispettive strutture differenziabili. 

    $ F: {M} \longrightarrow {N}$ si dice \emph{differenziabile} (o $C^{\infty}$) in $ p \in M $ se esistono una carta $ (U, \varphi) \in \mathcal{A} $ e una carta $ (V, \psi) \in\mathcal{B} $ tali che \begin{itemize}
        \item $ F(U) \subseteq V $; 
        \item $ \psi \circ F\circ \varphi^{-1} : \varphi(U) \longrightarrow \psi(V) $ sia di classe $ C^{\infty} $ in un intorno di $ \varphi(p) $.
    \end{itemize}

    Se $ F $ è differenziabile in ogni punto di $ M $, si dirà che $ F $ è \emph{differenziabile} (o $C^{\infty}$).

    Un'applicazione differenziabile, biiettiva e con inversa differenziabile è detta \emph{diffeomorfismo}.
}
\osservazione{
    La buona definizione di questa definizione segue dalla successiva proposizione.
}
\proposizione{jskdjhsfkjhcakjhdvcbloeiruhfcalfs}{
    Sia $F: M \longrightarrow N$ un'applicazione fra varietà, differenziabile in $p \in M$. Allora per ogni carta $(\tilde{U}, \tilde{\varphi})$ in $p$ e ogni carta $(\tilde{V}, \tilde{\psi})$ in $F(p)$ la composizione $\tilde{\psi} \circ F \circ \tilde{\varphi}^{-1}$ è differenziabile in $\tilde{\varphi}(p)$.
}
\proposizione{jnkjnkjnjdjdj}{
    Se $ M, N, P $ sono varietà differenziabili e \begin{equation*}
        F: M \longrightarrow N,\quad G: N \longrightarrow P 
    \end{equation*}sono mappe differenziabili, allora $ G\circ F: M \longrightarrow P $ è differenziabile.
}

\subsection{Vettori e spazio tangente}

\definizione{
    Sia $ M $ una varietà e $ p \in M $. Si definisce una relazione di equivalenza su \begin{equation*}
        \set{(U,f): U\text{ intorno aperto di }p, f: U \longrightarrow \R\text{ differenziabile}}
    \end{equation*}dove $ (U,f) \sim (V,g) $ se e solo se esiste $ W \subseteq U\cap V $ tale che $ \restriction{f}{W} = \restriction{g}{W} $.

    L'insieme delle classi di equivalenza è denotato con $ C^{\infty}(p) $.
}
\paragrafo{Lemma}{%
    Sia $ M $ una varietà e $ p \in M $. Siano $ \bm{f},\bm{g} \in C^{\infty}(p) $ due classi di equivalenza, e siano $ (U_1,f_1),(U_2,f_2) $ rappresentanti di $ \bm{f} $ e $ (V_1,g_1),(V_2,g_2) $ rappresentanti di $ \bm{g} $. Allora \begin{enumerate}
        \item $ (U_1\cap V_1, f_1+g_1)\sim (U_2\cap V_2, f_2+g_2) $;
        \item $ (U_1\cap V_1, f_1 \cdot g_1)\sim (U_2\cap V_2, f_2 \cdot g_2) $;
        \item $ (U_1, \lambda f_1) \sim (U_2, \lambda f_2)$ per ogni $\lambda \in \R$;
        \item $ f_1(p)=f_2(p) $.
    \end{enumerate}
}{}{}
\osservazione{
    Per $ \bm{f} \in C^{\infty}(p) $ rappresentato da $ (U,f) $ ha senso parlare di $ \bm{f}(p) \coloneqq f(p) $.
}
\definizione{
    Si dota $ C^{\infty}(p) $ di somma, prodotto e prodotto per scalari. Siano $ \bm{f},\bm{g} \in C^{\infty}(p) $, con $ (U,f)$ rappresentante di $ \bm{f} $, $ (V,g) $ rappresentante di $ \bm{g} $. \begin{itemize}
        \item $ \bm{f} + \bm{g} \in C^{\infty}(p) $ è rappresentata da $ (U\cap V, f+g) $;
        \item $ \bm{f} \cdot \bm{g} \in C^{\infty}(p) $ è rappresentata da $ (U\cap V, f \cdot g) $;
        \item $ \lambda \bm{f} \in C^{\infty}(p)$ rappresentato da $ (U,\lambda f)$ per ogni $\lambda \in \R$. 
    \end{itemize}
}
\definizione{
    Sia $ M $ una varietà e sia $ p \in M $. Una \emph{derivazione} su $ M $ in $ p $ è una applicazione $ \R $-lineare $ X: C^{\infty}(p) \longrightarrow \R $ tale che \begin{equation*}
        X(\bm{f}\bm{g}) = X(\bm{f}) \bm{g}(p) + \bm{f}(p)X(\bm{g})
    \end{equation*}per ogni $ \bm{f},\bm{g} \in C^{\infty}(p) $.
}
\definizione{
    L'insieme delle derivazioni di $ M $ in $ p $ è detto \emph{spazio tangente} ad $ M $ in $ p $, ed è denotato con $ T_{p}M  $. I suoi elementi si dicono \emph{vettori tangenti} ad $ M $ in $ p $.
}
\osservazione{
    $ T_{p}M  $ è uno spazio vettoriale reale.
}
\definizione{
    Sia $ M $ una varietà differenziabile, $ p \in M $ e $ \varphi = (u_1,\dots,u_{n}) $ una carta. Si definiscono le derivazioni $ \pd{}{u_{i}} $ che agiscono su $ \bm{f} \in C^{\infty}(p) $ come \begin{equation*}
        \restriction{\dpd{\bm{f}}{u_{i} }}{p} \coloneqq \restriction{\dpd{f \circ \varphi^{-1}}{x_{i} }}{\varphi(p)}
    \end{equation*}dove $ (U,f) $ è un rappresentante di $\bm{f}$.
}
\teorema{ksjnksjjdjdjcskjnsj}{
    Sia $ M $ una varietà differenziabile, $ p \in M $ e $ \varphi = (u_1,\dots,u_{n}) $ una carta. L'insieme \begin{equation*}
        \set{\restriction{\dpd{}{u_1}}{p},\dots,\restriction{\dpd{}{u_{n} }}{p}}
    \end{equation*}è una base dello spazio vettoriale $ T_{p}M  $.
}
\esempio{
    Sia $ \gamma : (- \varepsilon, \varepsilon) \longrightarrow M $ una curva di classe $ C^{\infty} $ tale che $ \gamma(0)=p \in M $. Il vettore $ \gamma'(0) $ tangente ad $ M $ in $ p $ agisce su $ \bm{f} \in C^{ \infty}(p)$ come \begin{equation*}
        \gamma'(0)(\bm{f}) \coloneqq \dod{(f\circ \gamma)}{t}(0)
    \end{equation*}dove $ (U,f) $ è un rappresentante di $ \bm{f} $. Se inoltre $ \varphi = (u_1,\dots,u_{n} ) $ è una carta definita in un intorno di $ p $, si ha che, in coordinate, scrivendo $ \varphi\circ \gamma = (\gamma^{1},\dots,\gamma^{n}) $\begin{equation*}
        \gamma'(0) = \sum_{j=1}^{n} (\gamma^{j})'(0) \restriction{\dpd{}{u_{j}}}{p}. 
    \end{equation*}
}
\subsection{Campi vettoriali e flusso}

\definizione{
    Un \emph{campo vettoriale} di classe $ C^{l} $ su $ M $ è una mappa che a ciascun punto $ p \in M $ assegna un vettore $ X_{p} \in T_{p}M $. In particolare, se in coordinate \begin{equation*}
        X = \sum_{j=1}^{n} X_{j}(p)\restriction{\dpd{}{u_{j}}}{p}
    \end{equation*}tutte le $ X_{j} $ devono essere funzioni $ C^{k} $.
}
\definizione{
    Sia $ X $ un campo vettoriale su $ M $ varietà differenziabile, e $ p \in M $. Una curva $ \gamma: (- \varepsilon, \varepsilon) \longrightarrow M $ tale che $ \gamma(0)= 0 $ e \begin{equation*}
        \gamma'(t) = X\left(\gamma(t)\right)
    \end{equation*}per ogni $ t \in (- \varepsilon, \varepsilon) $ è detta \emph{curva integrale} di $ X $ uscente da $ p $.
}
\osservazione{
    Per il teorema di Cauchy-Lipschitz sull'unicità della soluzione di un problema di Cauchy, se $ X $ è un campo vettoriale almeno $ C^{1} $ tale curva integrale esiste sempre su un intervallo massimale $ I $, contente lo zero.
}
\definizione{
    Sia $ X $ un campo vettoriale su una varietà $ M $. Si definisce il \emph{flusso di $ X $ al tempo $ t $} la mappa \begin{align*}
    \Theta_{t}:M &\longrightarrow M \\
    p &\longmapsto \gamma_{p}(t) 
    \end{align*}dove $ \gamma_{p}  $ è la curva integrale di $ X $ uscente da $ p $.
}